\documentclass[10pt,a4paper]{article}
\usepackage[top=2.0cm,bottom=2.0cm,left=2.0cm,right=2.0cm]{geometry}
\usepackage{fontspec}
\setmainfont{Times New Roman}
\setmonofont{Consolas}
\usepackage{amsmath,amssymb}
\usepackage{booktabs}
\usepackage{longtable}
\usepackage{tabularx}
\usepackage{array}
\usepackage{calc}
\usepackage{graphicx}
\usepackage{xcolor}
\usepackage{fancyhdr}
\usepackage{fancyvrb}
\usepackage{tcolorbox}
\usepackage{url}
\usepackage[colorlinks=true,linkcolor=blue!75!black,urlcolor=blue!75!black,citecolor=blue!75!black]{hyperref}

% Colors
\definecolor{primary}{RGB}{0,51,102}
\definecolor{accent}{RGB}{30,100,180}

% Header & Footer
\pagestyle{fancy}
\fancyhf{}
\fancyhead[L]{\small\itshape Explaining Neural Networks with Fuzzy Logic}
\fancyhead[R]{\small\itshape Technical Explainer \& User Guide}
\fancyfoot[C]{\small\thepage}
\renewcommand{\headrulewidth}{0.4pt}

% Disable automatic section numbering since markdown already has numbers
\setcounter{secnumdepth}{0}

% Section styling
\makeatletter
\renewcommand\section{\@startsection {section}{1}{\z@}%
                                   {-2.8ex \@plus -0.8ex \@minus -.2ex}%
                                   {1.4ex \@plus.2ex}%
                                   {\normalfont\Large\bfseries\color{primary}}}
\renewcommand\subsection{\@startsection{subsection}{2}{\z@}%
                                     {-2.4ex\@plus -0.6ex \@minus -.2ex}%
                                     {1.0ex \@plus .2ex}%
                                     {\normalfont\large\bfseries\color{primary}}}
\renewcommand\subsubsection{\@startsection{subsubsection}{3}{\z@}%
                                     {-2.0ex\@plus -0.5ex \@minus -.2ex}%
                                     {0.8ex \@plus .2ex}%
                                     {\normalfont\normalsize\bfseries\color{accent}}}
\makeatother

% Pandoc compatibility macros
\providecommand{\tightlist}{\setlength{\itemsep}{1pt}\setlength{\parskip}{1pt}}
\newcommand{\VerbBar}{|}
\newcommand{\VERB}{\Verb[commandchars=\\\{\}]}
\DefineVerbatimEnvironment{Highlighting}{Verbatim}{commandchars=\\\{\},fontsize=\small}
\newenvironment{Shaded}{\begin{quote}\small}{\end{quote}}

% Tokens
\newcommand{\AlertTok}[1]{\textcolor[rgb]{1.00,0.00,0.00}{\textbf{#1}}}
\newcommand{\AnnotationTok}[1]{\textcolor[rgb]{0.38,0.63,0.69}{\textbf{\textit{#1}}}}
\newcommand{\AttributeTok}[1]{\textcolor[rgb]{0.49,0.56,0.16}{#1}}
\newcommand{\BaseNTok}[1]{\textcolor[rgb]{0.25,0.63,0.44}{#1}}
\newcommand{\BuiltInTok}[1]{\textcolor[rgb]{0.00,0.50,0.00}{#1}}
\newcommand{\CharTok}[1]{\textcolor[rgb]{0.25,0.44,0.63}{#1}}
\newcommand{\CommentTok}[1]{\textcolor[rgb]{0.38,0.63,0.69}{\textit{#1}}}
\newcommand{\CommentVarTok}[1]{\textcolor[rgb]{0.38,0.63,0.69}{\textbf{\textit{#1}}}}
\newcommand{\ConstantTok}[1]{\textcolor[rgb]{0.53,0.00,0.00}{#1}}
\newcommand{\ControlFlowTok}[1]{\textcolor[rgb]{0.00,0.44,0.13}{\textbf{#1}}}
\newcommand{\DataTypeTok}[1]{\textcolor[rgb]{0.56,0.13,0.00}{#1}}
\newcommand{\DecValTok}[1]{\textcolor[rgb]{0.25,0.63,0.44}{#1}}
\newcommand{\DocumentationTok}[1]{\textcolor[rgb]{0.73,0.13,0.13}{\textit{#1}}}
\newcommand{\ErrorTok}[1]{\textcolor[rgb]{1.00,0.00,0.00}{\textbf{#1}}}
\newcommand{\ExtensionTok}[1]{#1}
\newcommand{\FloatTok}[1]{\textcolor[rgb]{0.25,0.63,0.44}{#1}}
\newcommand{\FunctionTok}[1]{\textcolor[rgb]{0.02,0.16,0.49}{#1}}
\newcommand{\ImportTok}[1]{\textcolor[rgb]{0.00,0.50,0.00}{\textbf{#1}}}
\newcommand{\InformationTok}[1]{\textcolor[rgb]{0.38,0.63,0.69}{\textbf{\textit{#1}}}}
\newcommand{\KeywordTok}[1]{\textcolor[rgb]{0.00,0.44,0.13}{\textbf{#1}}}
\newcommand{\NormalTok}[1]{#1}
\newcommand{\OperatorTok}[1]{\textcolor[rgb]{0.40,0.40,0.40}{#1}}
\newcommand{\OtherTok}[1]{\textcolor[rgb]{0.00,0.44,0.13}{#1}}
\newcommand{\PreprocessorTok}[1]{\textcolor[rgb]{0.74,0.48,0.00}{#1}}
\newcommand{\RegionMarkerTok}[1]{#1}
\newcommand{\SpecialCharTok}[1]{\textcolor[rgb]{0.25,0.44,0.63}{#1}}
\newcommand{\SpecialStringTok}[1]{\textcolor[rgb]{0.73,0.40,0.53}{#1}}
\newcommand{\StringTok}[1]{\textcolor[rgb]{0.25,0.44,0.63}{#1}}
\newcommand{\VariableTok}[1]{\textcolor[rgb]{0.10,0.09,0.49}{#1}}
\newcommand{\VerbatimStringTok}[1]{\textcolor[rgb]{0.25,0.44,0.63}{#1}}
\newcommand{\WarningTok}[1]{\textcolor[rgb]{0.38,0.63,0.69}{\textbf{\textit{#1}}}}

\setlength{\parindent}{0pt}
\setlength{\parskip}{4pt plus 1pt minus 1pt}

\begin{document}

\begin{center}
    \includegraphics[width=2.0cm,keepaspectratio]{report/figures/pec_logo.png}\\[0.15cm]
    {\fontsize{13}{16}\selectfont \bfseries PUNJAB ENGINEERING COLLEGE}\\[0.06cm]
    {\fontsize{11}{14}\selectfont \bfseries (DEEMED TO BE UNIVERSITY), CHANDIGARH}\\[0.06cm]
    {\fontsize{10}{13}\selectfont \bfseries Department of Electrical Engineering}\\[0.20cm]
    \rule{\linewidth}{1.0pt}\\[0.25cm]
    {\fontsize{14}{17}\selectfont \bfseries\color{primary} EXPLAINING NEURAL NETWORK DECISIONS IN\\ DIABETES DIAGNOSIS USING FUZZY LOGIC RULES}\\[0.15cm]
    {\fontsize{11}{14}\selectfont \bfseries\color{accent} Complete Technical Explainer \& Web Dashboard User Guide}\\[0.20cm]
    \rule{\linewidth}{0.6pt}\\[0.20cm]
\end{center}

\begin{tcolorbox}[colback=blue!3!white, colframe=primary, arc=2.5mm, boxrule=0.8pt, left=8pt, right=8pt, top=6pt, bottom=6pt]
\small
\textbf{Course:} Fuzzy Logic and Neural Networks (EEO1003) \hfill \textbf{Academic Year:} 2026--2027 \\
\textbf{Department:} Department of Electrical Engineering, Punjab Engineering College (Deemed to be University) \\
\textbf{Mentors:} Dr.~Rintu Khanna (Professor) \& Dr.~Hemant Kumar (Assistant Professor) \\
\textbf{Authors:} Rakshit (23106017), Aarav (23106023), Aryan Garg (23106024), Amish Bansal (23103087) \\
\textbf{Application Repository:} \texttt{d:\textbackslash{}Fuzzy\_Project\textbackslash{}Fuzzy\_Project} \hfill \textbf{Web Dashboard:} \url{http://127.0.0.1:8000}
\end{tcolorbox}
\vspace{0.3cm}








\subsection{Table of Contents}\label{table-of-contents}

\begin{enumerate}
\def\labelenumi{\arabic{enumi}.}
\tightlist
\item
  \hyperref[1-executive-summary-what-is-glass-box-insight]{Executive
  Summary: What is Glass-Box Insight?}
\item
  \hyperref[2-the-core-problem-why-is-this-project-needed]{The Core
  Problem: Why is this Project Needed?}
\item
  \hyperref[3-the-dataset-clinical-context--preprocessing-architecture]{The
  Dataset: Clinical Context \& Preprocessing Architecture}
\item
  \hyperref[4-the-core-techniques-used-step-by-step-technical-breakdown]{The
  Core Techniques Used: Step-by-Step Technical Breakdown}
\item
  \hyperref[5-how-to-launch-and-run-the-project]{How to Launch and Run
  the Project}
\item
  \hyperref[6-complete-web-dashboard-walkthrough-what-every-element-means]{Complete
  Web Dashboard Walkthrough (What Every Element Means)}

  \begin{itemize}
  \tightlist
  \item
    \hyperref[topbar--global-controls]{Topbar \& Global Controls}
  \item
    \hyperref[view-1-train-settings-stepper--real-time-telemetry]{View
    1: Train (Settings, Stepper, \& Real-Time Telemetry)}
  \item
    \hyperref[view-2-explore-clinical-sliders-dual-gauges--traceability]{View
    2: Explore (Clinical Sliders, Dual Gauges, \& Traceability)}
  \item
    \hyperref[view-3-edit-rules-human-in-the-loop-counterfactual-simulation]{View
    3: Edit Rules (Human-in-the-Loop Counterfactual Simulation)}
  \item
    \hyperref[view-4-compare-runs-pareto-frontier--historical-benchmarks]{View
    4: Compare Runs (Pareto Frontier \& Historical Benchmarks)}
  \end{itemize}
\item
  \hyperref[7-frequently-asked-questions--professor-viva-defense-guide]{Frequently
  Asked Questions \& Professor Viva Defense Guide}
\end{enumerate}

\begin{center}\rule{0.5\linewidth}{0.5pt}\end{center}

\subsection{1. Executive Summary: What is Glass-Box
Insight?}\label{executive-summary-what-is-glass-box-insight}

\textbf{Glass-Box Insight} is an Explainable Artificial Intelligence
(XAI) system designed for clinical decision support. Its core philosophy
is summarized by its tagline:
\[\textbf{\textit{``A neural network decides. Fuzzy rules explain why.''}}\]

\subsubsection{The High-Level Concept}\label{the-high-level-concept}

\begin{itemize}
\tightlist
\item
  In modern medicine, deep learning models and multi-layer neural
  networks achieve high diagnostic accuracy, but they operate as
  \textbf{opaque ``black boxes''}. They output a single probability
  number (e.g., \emph{``Patient \#152 has an 84\% probability of
  diabetes''}), but they cannot tell the doctor \emph{why}.
\item
  \textbf{Glass-Box Insight} solves this problem by training a
  transparent \textbf{Takagi-Sugeno-Kang (TSK) Neuro-Fuzzy Inference
  System} (ANFIS) to act as a \textbf{pedagogical surrogate} that
  imitates the neural network.
\item
  Instead of uninterpretable weight matrices, our system extracts
  \textbf{4 human-readable IF-THEN rules} with smooth linguistic terms
  (\emph{Low}, \emph{Medium}, \emph{High}), such as:
  \[\textbf{IF } \text{Glucose is High } \textbf{AND } \text{BMI is High } \textbf{THEN } \text{Risk} = 0.85 \text{ (High Risk)}\]
\item
  The extracted rules achieve \textbf{92.2\% test fidelity} (and
  \textbf{98.3\% confident fidelity}) relative to the neural network,
  meaning a physician gets the high accuracy of neural computation
  alongside the complete safety, auditability, and trust of transparent
  fuzzy logic.
\end{itemize}

\begin{center}\rule{0.5\linewidth}{0.5pt}\end{center}

\subsection{2. The Core Problem: Why is this Project
Needed?}\label{the-core-problem-why-is-this-project-needed}

\subsubsection{2.1 The High Accuracy vs.~Dangerous Opacity
Dilemma}\label{the-high-accuracy-vs.-dangerous-opacity-dilemma}

Artificial Neural Networks (ANNs) excel at finding non-linear patterns
in patient records. However, in critical healthcare applications,
deploying an opaque model introduces severe risks: 1.
\textbf{Verification Necessity:} A doctor cannot prescribe aggressive
diabetes medication or order expensive biopsies simply because a
computer said so without knowing which clinical biomarkers drove the
diagnosis. 2. \textbf{Confounding Artifacts:} Black boxes can exploit
dataset biases, laboratory recording errors, or spurious correlations
while maintaining high test accuracy. 3. \textbf{Medical Ethics \&
Doctor-Patient Communication:} Physicians are legally and ethically
obligated to explain diagnostic decisions in plain clinical terms to
patients. 4. \textbf{Regulatory Mandates:} - \textbf{EU AI Act:}
Formally classifies medical AI diagnostic tools as \emph{High-Risk},
mandating human interpretability and traceable decision oversight. -
\textbf{GDPR (Articles 13--22):} Guarantees patients the legal
\emph{``Right to Explanation''} for automated decisions. - \textbf{U.S.
FDA SaMD Guidelines:} Mandates full audit trails and decision
traceability for Software as a Medical Device.

\subsubsection{2.2 Why Existing Explainability Tools (LIME, SHAP,
Decision Trees) Fall
Short}\label{why-existing-explainability-tools-lime-shap-decision-trees-fall-short}

Prior approaches attempted to explain neural networks, but each has
major clinical flaws: - \textbf{LIME (Local Interpretable Model-agnostic
Explanations):} Perturbs input values randomly and fits local linear
models. Because it relies on stochastic sampling, explanations fluctuate
wildly across runs for the exact same patient. Furthermore, it only
explains a single patient and cannot provide an auditable hospital-wide
rule policy. - \textbf{SHAP (Shapley Additive Explanations):} Provides
local feature attribution bars, but assumes feature independence (which
fails in human physiology where Glucose and Insulin are tightly
correlated) and cannot provide global IF-THEN clinical rules. -
\textbf{Decision Tree Surrogates (TREPAN):} Uses rigid step-function
thresholds (e.g., \(\text{Glucose} \le 127.5\)). A patient with Glucose
= 127.4 mg/dL is classified as completely healthy, while 127.6 mg/dL
triggers a diabetic branch---an unrealistic, discontinuous cliff that
contradicts continuous human physiology.

\subsubsection{2.3 The Glass-Box Solution}\label{the-glass-box-solution}

By bridging \textbf{Fuzzy Set Theory} (continuous, graded membership
without cliffs) and \textbf{Neural Networks} (adaptive gradient
learning), Glass-Box Insight provides: - Smooth, realistic physiological
boundaries. - Global, auditable IF-THEN rule policies. - Full whiteboard
arithmetic traceability (every calculation can be verified by hand in 30
seconds). - Active detection of out-of-distribution / atypical patients
to prevent silent AI hallucinations.

\begin{center}\rule{0.5\linewidth}{0.5pt}\end{center}

\subsection{3. The Dataset: Clinical Context \& Preprocessing
Architecture}\label{the-dataset-clinical-context-preprocessing-architecture}

The system is evaluated on the benchmark \textbf{Pima Indians Diabetes
Database (PIDD)}, curated by the National Institute of Diabetes and
Digestive and Kidney Diseases (NIDDK) and hosted by the UCI Machine
Learning Repository.

\subsubsection{3.1 Clinical Features and Physiological
Axes}\label{clinical-features-and-physiological-axes}

The dataset contains records for \(N = 768\) female patients of Pima
Indian heritage aged 21 and older, categorized across three
physiological metabolic axes:

\begin{longtable}[]{@{}
  >{\raggedright\arraybackslash}p{(\linewidth - 6\tabcolsep) * \real{0.2353}}
  >{\raggedright\arraybackslash}p{(\linewidth - 6\tabcolsep) * \real{0.2353}}
  >{\centering\arraybackslash}p{(\linewidth - 6\tabcolsep) * \real{0.2941}}
  >{\raggedright\arraybackslash}p{(\linewidth - 6\tabcolsep) * \real{0.2353}}@{}}
\toprule\noalign{}
\begin{minipage}[b]{\linewidth}\raggedright
Physiological Axis
\end{minipage} & \begin{minipage}[b]{\linewidth}\raggedright
Feature Name
\end{minipage} & \begin{minipage}[b]{\linewidth}\centering
Physical Units
\end{minipage} & \begin{minipage}[b]{\linewidth}\raggedright
Clinical / Pathological Meaning
\end{minipage} \\
\midrule\noalign{}
\endhead
\bottomrule\noalign{}
\endlastfoot
\textbf{Glycemic Control} & \texttt{Glucose} & mg/dL & 2-hour plasma
glucose concentration in an oral glucose tolerance test (OGTT). Primary
diagnostic criterion for diabetes. \\
& \texttt{Insulin} & \(\mu\text{U/mL}\) & 2-hour post-prandial serum
insulin level. Reflects pancreatic beta-cell secretion capacity. \\
\textbf{Adiposity \& Body Mass} & \texttt{BMI} & \(\text{kg/m}^2\) &
Body Mass Index (weight in kg / \(\text{height}^2\) in m). Key indicator
of insulin resistance. \\
& \texttt{SkinThickness} & mm & Triceps skinfold thickness. Measures
subcutaneous adipose tissue deposits. \\
\textbf{Demographic \& Hereditary} & \texttt{Pregnancies} & Count &
Number of completed gestational cycles. Gestational diabetes stress
marker. \\
& \texttt{Age} & Years & Patient age at clinical examination. Beta-cell
function declines with age. \\
& \texttt{BloodPressure} & mm Hg & Diastolic blood pressure. Indicator
of vascular resistance and metabolic syndrome. \\
& \texttt{Pedigree} & Unitless & Diabetes Pedigree Function (DPF).
Scored genetic predisposition based on family history. \\
\textbf{Target Label} & \texttt{Outcome} & Binary & \(0 =\) Non-diabetic
(\(N=500, 65.1\%\)), \(1 =\) Diabetic (\(N=268, 34.9\%\)). \\
\end{longtable}

\subsubsection{3.2 The ``Zero-as-Missing'' Phenomenon \& Leakage-Free
Imputation}\label{the-zero-as-missing-phenomenon-leakage-free-imputation}

In the raw dataset, measurements that were omitted or unrecorded were
encoded as \texttt{0}. While \texttt{Pregnancies\ =\ 0} is biologically
valid, values such as \texttt{Glucose\ =\ 0},
\texttt{BloodPressure\ =\ 0}, or \texttt{BMI\ =\ 0} represent
\textbf{biological impossibilities}: - \texttt{Glucose}: 5 zero values
(\(0.65\%\)) - \texttt{BloodPressure}: 35 zero values (\(4.56\%\)) -
\texttt{BMI}: 11 zero values (\(1.43\%\)) - \texttt{SkinThickness}: 227
zero values (\(29.56\%\)) - \texttt{Insulin}: 374 zero values
(\(48.70\%\))

\textbf{Preprocessing Pipeline:} 1. \textbf{Stratified Split
(60/20/20):} Before any imputation or scaling, the dataset is
partitioned using fixed random seed 42 into: - \textbf{Training Set
(\(\mathcal{D}_{\text{train}}\)):} 460 patients (299 healthy, 161
diabetic) - \textbf{Validation Set (\(\mathcal{D}_{\text{val}}\)):} 154
patients (100 healthy, 54 diabetic) - \textbf{Test Set
(\(\mathcal{D}_{\text{test}}\)):} 154 patients (101 healthy, 53
diabetic) 2. \textbf{Leakage-Free Imputation:} Medians are calculated
\textbf{strictly on positive values from the training partition only}: -
\texttt{Glucose}: 117.0 mg/dL - \texttt{BloodPressure}: 72.0 mm Hg -
\texttt{SkinThickness}: 29.0 mm - \texttt{Insulin}: 125.0
\(\mu\text{U/mL}\) - \texttt{BMI}: 32.3 \(\text{kg/m}^2\) Validation and
test missing values are imputed using these exact training medians,
preventing data leakage. 3. \textbf{Dual Feature Scaling:} -
\textbf{Standardization (Z-score):}
\(\tilde{x} = (x - \mu_{\text{train}}) / \sigma_{\text{train}}\) for the
MLP neural network to ensure smooth gradient descent. - \textbf{Min-Max
Normalization:}
\(\tilde{x} = (x - x_{\text{min}}) / (x_{\text{max}} - x_{\text{min}}) \in [0, 1]\)
for the ANFIS neuro-fuzzy system to establish bounded, interpretable
linguistic universes of discourse. 4. \textbf{Synthetic Oracle Sampling
(\(N_{\text{synth}} = 2{,}000\)):} We sample 2,000 synthetic patient
vectors uniformly within the clinical feature hypercube and query the
trained neural network to label them. This populates sparse boundary
regions, giving the fuzzy surrogate rich decision-boundary information.

\begin{center}\rule{0.5\linewidth}{0.5pt}\end{center}

\subsection{4. The Core Techniques Used: Step-by-Step Technical
Breakdown}\label{the-core-techniques-used-step-by-step-technical-breakdown}

\begin{verbatim}
Raw Patient Records
       |
       ▼
[Stage 1: Train Neural Network (MLP 8-8-1)] --► Produces Black-Box Predictions (ROC-AUC: 0.797)
       |
       ▼
[Stage 2: Permutation Feature Importance]  --► Isolates Top 4 Drivers: Glucose, BMI, Age, Pregnancies
       |
       ▼
[Stage 3: Fuzzy C-Means (FCM) Clustering] --► Discovers 8 Patient Prototypes (Prevents 3^4=81 explosion)
       |
       ▼
[Stage 4: Shouldered Gaussian Memberships]--► Low, Medium, High terms with flat saturation shoulders
       |
       ▼
[Stage 5: 5-Layer TSK ANFIS Architecture]  --► Fuzzification → T-Norm → Normalize → Consequent → Sum
       |
       ▼
[Stage 6: Hybrid Learning Optimization]   --► Ridge Least Squares (Consequents) + Adam (Premises)
       |
       ▼
[Stage 7: Semantic Integrity Constraints] --► Enforces Low < Med < High & Drift bounds (Anti-Collapse)
       |
       ▼
[Stage 8: Greedy Backward Rule Pruning]   --► Prunes conditions & rules (8 rules → 4 compact rules)
       |
       ▼
[Stage 9: Clinical Deployment & Telemetry]--► 92.2% Test Fidelity, Sub-5ms Live Interactive Web Dashboard
\end{verbatim}

\subsubsection{4.1 Stage 1: The Black-Box Neural Network
(StreamingMLP)}\label{stage-1-the-black-box-neural-network-streamingmlp}

\begin{itemize}
\tightlist
\item
  \textbf{Architecture:} Feedforward Multi-Layer Perceptron (8 input
  features \(\to\) 8 hidden neurons with ReLU activation \(\to\) 1
  output neuron with Sigmoid activation).
\item
  \textbf{Complexity:} Only 81 trainable parameters
  (\((8 \times 8 + 8) + (8 \times 1 + 1) = 81\)), avoiding severe
  overfitting on small clinical cohorts.
\item
  \textbf{Optimization:} Adam optimizer (\(\beta_1=0.9, \beta_2=0.999\),
  learning rate \(\eta=0.01\)), regularized binary cross-entropy loss
  with L2 penalty (\(\alpha=1.0\)).
\item
  \textbf{Performance:} 70.1\% test accuracy, 0.797 ROC-AUC, 5-fold CV
  score of \(0.826 \pm 0.017\).
\end{itemize}

\subsubsection{4.2 Stage 2: Permutation Feature Importance
(PFI)}\label{stage-2-permutation-feature-importance-pfi}

If we fed all 8 features into a fuzzy system with 3 linguistic terms
each, grid partitioning would generate
\(3^8 = \mathbf{6{,}561 \text{ rules}}\)---an impossible cognitive
burden. To reduce dimensionality principledly, we measure how much the
neural network's cross-entropy loss degrades when each feature column is
randomly shuffled: 1. \textbf{Glucose:} 43.8\% importance share (Loss
increase +0.231) --- \textbf{Selected (Rank 1)} 2. \textbf{BMI:} 20.8\%
importance share (Loss increase +0.110) --- \textbf{Selected (Rank 2)}
3. \textbf{Age:} 9.7\% importance share (Loss increase +0.051) ---
\textbf{Selected (Rank 3)} 4. \textbf{Pregnancies:} 8.6\% importance
share (Loss increase +0.045) --- \textbf{Selected (Rank 4)} 5.
\emph{Non-selected features:} BloodPressure (5.9\%), Insulin (4.8\%),
SkinThickness (3.4\%), Pedigree (3.0\%). Selecting the top 4 salient
features captures \textbf{82.9\% of the network's predictive mass},
while shrinking the input dimensionality from 8 down to 4.

\subsubsection{4.3 Stage 3: Data-Driven Fuzzy C-Means (FCM)
Clustering}\label{stage-3-data-driven-fuzzy-c-means-fcm-clustering}

Even with 4 features, full grid partitioning produces
\(3^4 = \mathbf{81 \text{ rules}}\). To solve this, we run \textbf{Fuzzy
C-Means clustering} (Bezdek, 1981) with fuzziness exponent \(m=2.0\) on
the 4-dimensional normalized patient data. FCM discovers \(C=8\)
multi-dimensional cluster centroids (patient prototypes) that reflect
actual clinical patterns, directly initializing \textbf{8 candidate
rules}.

\subsubsection{4.4 Stage 4: Shouldered Gaussian Membership
Functions}\label{stage-4-shouldered-gaussian-membership-functions}

Standard Gaussian curves decay to zero at extremes. For a patient with
extreme glucose (\(\ge 200\) mg/dL), a standard Gaussian would output
near-zero membership, leaving the patient with no firing rules! We
implement \textbf{monotonic flat boundary shoulders}: -
\textbf{\texttt{Low}:} Gaussian curve for \(x \ge c_{\text{Low}}\), but
remains \textbf{flat at 1.0} for all \(x \le c_{\text{Low}}\). -
\textbf{\texttt{Medium}:} Standard bell-shaped Gaussian centered at
\(c_{\text{Med}}\). - \textbf{\texttt{High}:} Gaussian curve for
\(x \le c_{\text{High}}\), but remains \textbf{flat at 1.0} for all
\(x \ge c_{\text{High}}\). This guarantees that severely ill patients
always receive full membership in \emph{High}, completely eliminating
unhandled edge cases.

\subsubsection{4.5 Stage 5: Five-Layer Takagi-Sugeno-Kang (TSK) ANFIS
Architecture}\label{stage-5-five-layer-takagi-sugeno-kang-tsk-anfis-architecture}

The extracted system is executed through 5 transparent feedforward
layers: - \textbf{Layer 1 (Fuzzification):} Computes continuous
membership grades \(\mu_{jt}(x_j) \in [0, 1]\) for each feature in
\emph{Low}, \emph{Medium}, or \emph{High}. - \textbf{Layer 2 (Rule
Firing via Product T-Norm):} Multiplies memberships across the active
conditions of rule \(r\):
\[w_r = \prod_{j \in \text{Cond}(r)} \mu_{j, A_{rj}}(x_j)\] We use the
differentiable Product operator (\(\times\)) rather than the Min
(\(\min\)) operator, ensuring smooth gradient propagation. -
\textbf{Layer 3 (Normalization):} Calculates the relative firing weight
of each rule:
\[\bar{w}_r = \frac{w_r}{\sum_{k=1}^R w_k}, \quad \text{such that } \sum_{r=1}^R \bar{w}_r = 1.0\]
- \textbf{Layer 4 (Rule Consequents):} In zero-order TSK, each rule
outputs a fixed risk scalar \(p_r \in [0, 1]\). - \textbf{Layer 5
(Defuzzified Output):} Calculates the final patient diabetes risk score
as a weighted sum: \[y = \sum_{r=1}^R \bar{w}_r \cdot p_r\]

\subsubsection{4.6 Stage 6: Hybrid Learning
Optimization}\label{stage-6-hybrid-learning-optimization}

ANFIS updates its parameters through a fast, stable two-step hybrid
algorithm: 1. \textbf{Forward Pass (Least Squares for Consequents):}
With Gaussian membership curves held fixed, the rule output risk weights
\(\mathbf{p} = [p_1, \dots, p_R]^T\) are solved in \textbf{a single
mathematical step} using Ridge Regularized Least Squares (RLS):
\[\mathbf{p}^* = (\bar{\mathbf{W}}^T \bar{\mathbf{W}} + \lambda \mathbf{I})^{-1} \bar{\mathbf{W}}^T \mathbf{t}\]
where \(\mathbf{t}\) contains the continuous probabilities output by the
neural network. 2. \textbf{Backward Pass (Gradient Descent for
Premises):} With consequents \(\mathbf{p}\) fixed, prediction errors
\((y_i - t_i)\) are backpropagated to update the Gaussian centers
(\(c\)) and widths (\(\sigma\)) using the Adam optimizer
(\(\eta=0.01\)).

\subsubsection{4.7 Stage 7: Semantic Integrity Preservation
Constraints}\label{stage-7-semantic-integrity-preservation-constraints}

Unconstrained gradient descent can cause \textbf{semantic
collapse}---where the \emph{Low} center drifts higher than the
\emph{High} center, or curves shrink into needle-thin spikes. After
every training iteration, three geometric constraints are enforced: 1.
\textbf{Strict Term Ordering:}
\(c_{\text{Low}} < c_{\text{Med}} < c_{\text{High}}\) (enforced with a
minimum gap of \(0.05\)). 2. \textbf{Bounded Anchor Drift:} Centers
cannot drift more than \(\pm 0.10\) from their original data-derived
cluster centers (\(|c - c^{(0)}| \le 0.10\)). 3. \textbf{Minimum Width
Floor:} Gaussian widths are bounded below by \(\sigma \ge 0.05\) to
prevent delta-function spikes.

\subsubsection{4.8 Stage 8: Greedy Backward Rule
Pruning}\label{stage-8-greedy-backward-rule-pruning}

The initialized 8-rule system contains redundant conditions (e.g., if a
rule fires whenever Glucose is High regardless of Age, Age is
redundant). We run an automated \textbf{two-check greedy backward
elimination} algorithm: 1. Iteratively test deleting each condition one
by one. 2. An edit is accepted if and only if: - Holdout validation
fidelity drops by less than the user tolerance
(\(\Delta \mathcal{F}_{\text{val}} \le 0.01\)). - Absolute patient
probability error remains bounded (\(\text{MAE} \le 0.05\)). 3. If all
conditions in a rule become empty, the entire rule is pruned.
\textbf{Result:} The rule base drops from \textbf{8 rules and 32
conditions} down to \textbf{4 compact rules and 12 conditions}---a
\textbf{96\% complexity reduction} from the theoretical grid while
retaining \textbf{92.2\% test fidelity}.

\subsubsection{4.9 Stage 9: Safety Diagnostics \&
Traceability}\label{stage-9-safety-diagnostics-traceability}

\begin{itemize}
\tightlist
\item
  \textbf{Hand Traceability:} The exact calculation for any patient can
  be computed on a whiteboard in 30 seconds using standard arithmetic.
\item
  \textbf{Weak-Match / Out-of-Distribution Diagnostic Alert:} The
  maximum raw firing strength is tracked:
  \[\mathcal{S}_{\text{match}} = \max_{r=1 \dots R} w_r\] If an atypical
  patient arrives whose clinical profile does not match any known
  prototype (e.g., an elderly, severely malnourished patient with
  extreme glucose), \(\mathcal{S}_{\text{match}}\) drops below \(0.20\).
  The system immediately triggers an amber/red \textbf{``Weak Match: Out
  of Distribution''} alert, warning doctors that the automated system
  cannot reliably evaluate this patient---preventing silent AI errors.
\end{itemize}

\begin{center}\rule{0.5\linewidth}{0.5pt}\end{center}

\subsection{5. How to Launch and Run the
Project}\label{how-to-launch-and-run-the-project}

\subsubsection{5.1 System Requirements \&
Verification}\label{system-requirements-verification}

The project runs on standard Python 3.10+ (tested on Python 3.13) with
standard scientific computing packages: - \texttt{fastapi},
\texttt{uvicorn}, \texttt{starlette} (Web server \& REST API) -
\texttt{numpy}, \texttt{pandas}, \texttt{scipy} (Data handling \& linear
algebra) - \texttt{pytest} (Automated testing suite)

\subsubsection{5.2 Starting the Web
Server}\label{starting-the-web-server}

Open PowerShell or your command terminal and navigate to the project
root:

\begin{Shaded}
\begin{Highlighting}[]
\FunctionTok{cd}\NormalTok{ d}\OperatorTok{:}\NormalTok{\textbackslash{}Fuzzy\_Project\textbackslash{}Fuzzy\_Project}
\NormalTok{uvicorn webapp}\OperatorTok{.}\FunctionTok{server}\OperatorTok{:}\NormalTok{app }\OperatorTok{{-}{-}}\NormalTok{host }\DecValTok{127.0}\OperatorTok{.}\DecValTok{0.1} \OperatorTok{{-}{-}}\NormalTok{port }\DecValTok{8000}
\end{Highlighting}
\end{Shaded}

When successful, the terminal displays:

\begin{verbatim}
INFO:     Started server process
INFO:     Waiting for application startup.
INFO:     Application startup complete.
INFO:     Uvicorn running on http://127.0.0.1:8000 (Press CTRL+C to quit)
\end{verbatim}

\subsubsection{5.3 Opening the Dashboard}\label{opening-the-dashboard}

Open your browser and navigate to: \textbf{\url{http://127.0.0.1:8000}}

\subsubsection{5.4 Running the Automated Test
Suite}\label{running-the-automated-test-suite}

To verify the integrity of all algorithmic modules, execute pytest:

\begin{Shaded}
\begin{Highlighting}[]
\FunctionTok{cd}\NormalTok{ d}\OperatorTok{:}\NormalTok{\textbackslash{}Fuzzy\_Project\textbackslash{}Fuzzy\_Project}
\NormalTok{pytest tests}\OperatorTok{/}
\end{Highlighting}
\end{Shaded}

All \textbf{26 tests} will pass in \textasciitilde6 seconds: -
\texttt{test\_data.py}: Verifies stratified splits, median imputation,
and normalization. - \texttt{test\_blackbox.py}: Verifies MLP forward
pass and permutation feature importance. - \texttt{test\_fcm.py}:
Verifies Fuzzy C-Means convergence and membership constraints. -
\texttt{test\_anfis.py}: Verifies 5-layer forward propagation, Ridge
RLS, and hybrid learning. - \texttt{test\_metrics\_and\_rules.py}:
Verifies greedy pruning and fidelity metrics. -
\texttt{test\_webapp.py}: Verifies FastAPI REST endpoints and SSE
telemetry streams.

\begin{center}\rule{0.5\linewidth}{0.5pt}\end{center}

\subsection{6. Complete Web Dashboard Walkthrough (What Every Element
Means)}\label{complete-web-dashboard-walkthrough-what-every-element-means}

The web dashboard is divided into four main functional views accessible
via the top navigation bar.

\begin{verbatim}
+------------------------------------------------------------------------------+
|  GLASS-BOX INSIGHT     [1 · Train] [2 · Explore] [3 · Edit rules] [4 · Runs] |
+------------------------------------------------------------------------------+
\end{verbatim}

\begin{center}\rule{0.5\linewidth}{0.5pt}\end{center}

\subsubsection{Topbar \& Global Controls}\label{topbar-global-controls}

\begin{verbatim}
+------------------------------------------------------------------------------+
| [Logo] Glass-Box Insight                         [Run #1 · 4 rules]  [ ◐ ]   |
|        A neural network decides. Fuzzy rules explain why.                    |
+------------------------------------------------------------------------------+
\end{verbatim}

\begin{enumerate}
\def\labelenumi{\arabic{enumi}.}
\tightlist
\item
  \textbf{Brand \& Tagline:} Identifies the system and its pedagogical
  neuro-fuzzy surrogate role.
\item
  \textbf{Navigation Tabs:}

  \begin{itemize}
  \tightlist
  \item
    \texttt{1\ ·\ Train}: Configure models and watch the complete
    training pipeline run live.
  \item
    \texttt{2\ ·\ Explore}: Interactive clinical sliders and real-time
    dual-model patient diagnosis.
  \item
    \texttt{3\ ·\ Edit\ rules}: Live interactive rule editor for
    human-in-the-loop counterfactual simulation.
  \item
    \texttt{4\ ·\ Compare\ runs}: Pareto frontier scatter chart and
    historical run logs.
  \end{itemize}
\item
  \textbf{Run Pill (\texttt{\#run-pill}):}

  \begin{itemize}
  \tightlist
  \item
    Shows the status of the currently loaded model (e.g.,
    \texttt{No\ run\ yet} or
    \texttt{Run\ \#1\ ·\ 4\ rules\ ·\ 92.2\%\ fidelity}).
  \item
    Turns green on successful training completion or yellow while
    training is actively streaming.
  \end{itemize}
\item
  \textbf{Theme Toggle (\texttt{◐}):} Instantly toggles between
  high-contrast Dark Mode and clean Clinical Light Mode.
\end{enumerate}

\begin{center}\rule{0.5\linewidth}{0.5pt}\end{center}

\subsubsection{View 1: Train (Settings, Stepper, \& Real-Time
Telemetry)}\label{view-1-train-settings-stepper-real-time-telemetry}

This view allows you to configure hyper-parameters, initiate training,
and watch the system learn in real time via Server-Sent Events (SSE).

\begin{verbatim}
+---------------------------+--------------------------------------------------+
|  SETTINGS                 |  [Data]-►[Network]-►[Importance]-►[ANFIS]-►[Prune]|
|                           +-------------------------+------------------------+
|  Black Box Neural Network | ① Neural Network        | ② Feature Importance   |
|  - Hidden layers: 8       |   [Loss Decay Chart]    |   [PFI Bar Chart]      |
|  - L2 penalty α: 1        +-------------------------+------------------------+
|                           | ③ ANFIS Learns to Imitate Network                |
|  Fuzzy Explainer (ANFIS)  |   [Hybrid Loss Chart] | [Membership Curves Low/H]|
|  - Inputs: 4              +--------------------------------------------------+
|  - Terms: 3 (Low/Med/High)| ④ Rule Base: Proposed by FCM, Then Pruned        |
|  - Clusters: 8            |   [Active Rules List] | [Pruning Event Log]      |
|  - Target: The network    +--------------------------------------------------+
|                           | ⑤ Evaluation Results on 154 Test Patients        |
|  [ ▶ Train ]  [ Reset ]   |   [ 92.2% Fidelity ] [ 70.1% Acc ] [ 0.797 AUC ] |
+---------------------------+--------------------------------------------------+
\end{verbatim}

\paragraph{The Settings Form (Left
Sidebar)}\label{the-settings-form-left-sidebar}

\begin{itemize}
\tightlist
\item
  \textbf{Black box · neural network:}

  \begin{itemize}
  \tightlist
  \item
    \texttt{Hidden\ layers}: Selects the network depth
    (\texttt{8\ neurons} (default), \texttt{16\ neurons},
    \texttt{16\ →\ 8\ neurons}, or \texttt{32\ →\ 16\ neurons}).
  \item
    \texttt{L2\ penalty\ α}: Weight decay regularization
    (\(0.001, 0.01, 0.1, 1.0\)). Higher \(\alpha\) prevents overfitting.
  \item
    \texttt{Choose\ size\ and\ α\ by\ 5-fold\ cross-validation}:
    Checkbox that runs an automated grid search over architectures
    before training.
  \end{itemize}
\item
  \textbf{Fuzzy explainer · ANFIS:}

  \begin{itemize}
  \tightlist
  \item
    \texttt{Inputs\ used\ by\ the\ rules} (Slider 2 to 6): Controls how
    many salient features from PFI are passed to the fuzzy explainer
    (Default: 4).
  \item
    \texttt{Terms\ per\ input}: Choose between 2 linguistic terms
    (\emph{Low}, \emph{High}) or 3 terms (\emph{Low}, \emph{Medium},
    \emph{High}).
  \item
    \texttt{FCM\ clusters} (Slider 2 to 16): Upper bound on initial
    candidate rules (Default: 8).
  \item
    \texttt{THEN-part}: Choose between \textbf{Zero-order TSK} (Constant
    scalar, e.g., \(\text{Risk} = 0.85\)) or \textbf{First-order TSK}
    (Linear polynomial formula).
  \item
    \texttt{Training\ epochs} (Slider 20 to 600): Number of hybrid
    learning cycles (Default: 100).
  \item
    \texttt{Imitate}:

    \begin{itemize}
    \tightlist
    \item
      \texttt{The\ network\ (explainer)}: True surrogate mode. ANFIS
      imitates the neural network's continuous probabilities.
    \item
      \texttt{The\ true\ diagnosis\ (control)}: Trains the fuzzy system
      directly on ground-truth \(\{0, 1\}\) clinical labels (serves as
      an experimental control baseline).
    \end{itemize}
  \item
    \texttt{Keep\ terms\ meaningful\ (constraints)}: Enforces ordering
    (\(c_{\text{Low}} < c_{\text{Med}} < c_{\text{High}}\)) and prevents
    semantic collapse.
  \end{itemize}
\item
  \textbf{Data \& simplification:}

  \begin{itemize}
  \tightlist
  \item
    \texttt{Synthetic\ patients\ asked\ to\ the\ network} (0 to 4,000):
    Number of synthetic oracle neighborhood queries used to populate
    boundary regions (Default: 2,000).
  \item
    \texttt{Prune\ conditions\ and\ rules}: Enables the automated greedy
    backward pruning engine.
  \item
    \texttt{Allowed\ fidelity\ loss\ while\ pruning} (\(0\%\) to
    \(5\%\)): Maximum permitted validation fidelity drop per pruning
    step (Default: \(1\%\)).
  \item
    \texttt{Random\ seed}: Controls synthetic sampling and cluster
    initialization (Default: 42).
  \item
    \texttt{Slow\ motion}: Adds an artificial 25ms delay per epoch so
    users can visually watch the curves move live on projector screens.
  \end{itemize}
\item
  \textbf{Action Buttons:}

  \begin{itemize}
  \tightlist
  \item
    \texttt{▶\ Train}: Starts the asynchronous training pipeline.
  \item
    \texttt{Report\ settings}: Restores all parameters to the exact
    benchmark configuration reported in the academic paper.
  \end{itemize}
\end{itemize}

\paragraph{The 7-Stage Pipeline
Stepper}\label{the-7-stage-pipeline-stepper}

As training executes, the horizontal stepper highlights the active
algorithmic stage: 1. \texttt{Data}: Ingestion, stratified splitting,
and median imputation. 2. \texttt{Neural\ network}: Backpropagation
training of the MLP. 3. \texttt{Importance}: Permutation feature
importance shuffling. 4. \texttt{FCM\ rules}: Unsupervised prototype
clustering. 5. \texttt{ANFIS}: Hybrid optimization of premise curves and
consequent weights. 6. \texttt{Pruning}: Backward greedy elimination of
redundant conditions. 7. \texttt{Evaluation}: Final benchmarking across
the 154 holdout test patients.

\paragraph{Live Telemetry Dashboard
Cards}\label{live-telemetry-dashboard-cards}

\begin{itemize}
\tightlist
\item
  \textbf{Card ① Neural Network Learns:}

  \begin{itemize}
  \tightlist
  \item
    Displays the live training loss curve (Binary Cross-Entropy)
    decreasing across epochs.
  \item
    KPI badge shows final training loss and convergence epoch.
  \end{itemize}
\item
  \textbf{Card ② What the Network Relies On:}

  \begin{itemize}
  \tightlist
  \item
    Horizontal bar chart showing the permutation importance share of
    each clinical feature.
  \item
    The top features selected for the fuzzy explainer are highlighted in
    bright blue.
  \end{itemize}
\item
  \textbf{Card ③ ANFIS Learns to Imitate the Network:}

  \begin{itemize}
  \tightlist
  \item
    \textbf{Left:} Mean Squared Error (MSE) loss decay curve as ANFIS
    imitates the neural network's output probabilities.
  \item
    \textbf{Right (MF Grid):} Real-time animated Gaussian membership
    curves for Glucose, BMI, Age, and Pregnancies. Dashed curves
    represent the initial cluster positions; solid lines show the curves
    being fine-tuned live by gradient descent.
  \end{itemize}
\item
  \textbf{Card ④ Rule Base: Proposed by FCM, Then Pruned:}

  \begin{itemize}
  \tightlist
  \item
    \textbf{Left:} The live rule list. Rules that get completely pruned
    disappear from this list.
  \item
    \textbf{Right (Pruning Event Log):} Real-time text log of every
    accepted pruning decision (e.g.,
    \texttt{Pruned\ Age\ is\ High\ from\ Rule\ 3\ (Fidelity:\ 94.8\%\ -\textgreater{}\ 94.6\%)}).
  \end{itemize}
\item
  \textbf{Card ⑤ Evaluation Results:}

  \begin{itemize}
  \tightlist
  \item
    Unrolls once training finishes, presenting holdout test metrics:

    \begin{itemize}
    \tightlist
    \item
      \textbf{Fidelity:} Percentage of test patients where ANFIS and the
      Neural Network assign the identical diagnostic class (\(92.2\%\)).
    \item
      \textbf{NN Test Accuracy:} Diagnostic accuracy of the black box
      (\(70.1\%\)).
    \item
      \textbf{Surrogate Accuracy:} Diagnostic accuracy of the fuzzy
      explainer (\(72.1\%\)).
    \item
      \textbf{ROC-AUC:} Area under the ROC curve (\(0.797\)).
    \end{itemize}
  \item
    Action buttons allow instant transition to \textbf{``Explore
    patients →''} or \textbf{``Edit the rules →''}.
  \end{itemize}
\end{itemize}

\begin{center}\rule{0.5\linewidth}{0.5pt}\end{center}

\subsubsection{View 2: Explore (Clinical Sliders, Dual Gauges, \&
Traceability)}\label{view-2-explore-clinical-sliders-dual-gauges-traceability}

The Explore view is the primary clinical auditing workstation where a
physician can evaluate real patients or test hypothetical ``what-if''
scenarios.

\begin{verbatim}
+-----------------------+------------------------------------------------------+
|  PATIENT SELECTION    |  CONCURRENT DUAL VERDICT                             |
|  [ Patient #152    ▼] |   [ Black Box MLP ]         [ Glass-Box ANFIS ]      |
|                       |      83.7% Risk                 84.3% Risk           |
|  CLINICAL SLIDERS     |     (Diabetic)                 (Diabetic)            |
|  Glucose:    197 mg/dL|   [Badge: Perfect Agreement]  [Badge: Strong Match]  |
|  BMI:        34.7     +------------------------------------------------------+
|  Age:        31 yrs   |  WHY: THE RULES THAT FIRED                           |
|  Pregnancies: 2       |  Stacked Bar: [██████████████████████░░░░] 84.3% Risk|
|                       |  Rule 4 (Diabetic):  Weight 98.4% | Contribution 0.84|
|  [True Label: 1]      |  Rule 1 (Healthy):   Weight  1.6% | Contribution 0.00|
|                       +------------------------------------------------------+
|                       |  FUZZIFICATION BREAKDOWN                             |
|                       |  Glucose (197): Low 0.00 | Med 0.00 | High 1.00      |
|                       |  BMI (34.7):    Low 0.00 | Med 0.65 | High 0.35      |
+-----------------------+------------------------------------------------------+
\end{verbatim}

\paragraph{1. Patient Selector \& Input Sliders (Left
Sidebar)}\label{patient-selector-input-sliders-left-sidebar}

\begin{itemize}
\tightlist
\item
  \textbf{Start from Dropdown:} Allows selecting any of the 154 holdout
  test patients or predefined clinical prototypes:

  \begin{itemize}
  \tightlist
  \item
    \texttt{Diabetic\ Prototype}: Severe hyperglycemia and high BMI.
  \item
    \texttt{Healthy\ Prototype}: Low glucose, lean BMI, young age.
  \item
    \texttt{Boundary\ Ambiguity\ Case}: Borderline glucose (125 mg/dL)
    near clinical diagnostic thresholds.
  \item
    \texttt{Atypical\ Outlier\ Case}: Severe hyperglycemia in an elderly
    patient with lean BMI (triggers weak match).
  \end{itemize}
\item
  \textbf{Interactive Sliders:}

  \begin{itemize}
  \tightlist
  \item
    \texttt{Glucose} (\(0\) to \(200\) mg/dL, step \(1\))
  \item
    \texttt{BMI} (\(15.0\) to \(60.0 \text{ kg/m}^2\), step \(0.1\))
  \item
    \texttt{Age} (\(21\) to \(85\) years, step \(1\))
  \item
    \texttt{Pregnancies} (\(0\) to \(17\), step \(1\))
  \end{itemize}
\item
  Whenever a slider is dragged, an asynchronous call is dispatched to
  the Python backend; predictions, gauges, and rule activations update
  in \textbf{less than 5 milliseconds}.
\end{itemize}

\paragraph{2. Concurrent Dual Verdict
Gauges}\label{concurrent-dual-verdict-gauges}

\begin{itemize}
\tightlist
\item
  \textbf{Gauge 1 (Neural Network):} Displays the continuous probability
  \(P_{\text{NN}} \in [0\%, 100\%]\) produced by the multi-layer
  perceptron.
\item
  \textbf{Gauge 2 (Glass-Box Fuzzy):} Displays the exact defuzzified
  probability \(P_{\text{fuzzy}} \in [0\%, 100\%]\) produced by the
  ANFIS surrogate.
\item
  \textbf{Diagnostic Status Badges:}

  \begin{itemize}
  \tightlist
  \item
    \texttt{Agreement\ Badge}: Shows \texttt{Agreement} (both predict
    healthy or both predict diabetic) or \texttt{Disagreement} (if the
    models cross the 0.50 boundary).
  \item
    \texttt{Match-Strength\ Badge}: Shows \texttt{Strong\ Match}
    (\(\mathcal{S}_{\text{match}} \ge 0.50\)), \texttt{Moderate\ Match},
    or \texttt{Weak\ Match:\ Out\ of\ Distribution}
    (\(\mathcal{S}_{\text{match}} < 0.20\), highlighted in amber/red).
  \end{itemize}
\end{itemize}

\paragraph{3. Card ``Why: The Rules That
Fired''}\label{card-why-the-rules-that-fired}

This panel provides complete, transparent decision decomposition: -
\textbf{Stacked Probability Bar:} Visualizes total risk. Green segments
represent contributions from non-diabetic rules; red segments represent
contributions from diabetic rules. - \textbf{Fired Rules Breakdown
Table:} - \textbf{Rule Sentence:} Full linguistic text (e.g.,
\texttt{IF\ Glucose\ is\ High\ AND\ BMI\ is\ Medium\ THEN\ Risk\ =\ 0.85}).
- \textbf{Raw Firing Strength (\(w_r\)):} The product of membership
grades across active conditions. - \textbf{Normalized Weight
(\(\bar{w}_r\)):} The rule's percentage share of the overall clinical
decision (\(\sum \bar{w}_r = 100\%\)). - \textbf{Consequent Risk
(\(p_r\)):} The base risk score associated with this rule prototype. -
\textbf{Risk Contribution:} Exactly equals \(\bar{w}_r \times p_r\). The
sum of all contributions \textbf{adds up exactly to the final fuzzy risk
percentage}.

\paragraph{4. Card ``Fuzzification
Breakdown''}\label{card-fuzzification-breakdown}

Visualizes exactly how the patient's raw clinical measurements map into
linguistic terms: - Displays 4 plots (one for each feature) with
\emph{Low}, \emph{Medium}, and \emph{High} curves. - A vertical
indicator line marks the patient's current slider value. - Clear numeric
tags show the exact membership grade (e.g.,
\(\mu_{\text{High}}(\text{Glucose}) = 0.98\)).

\begin{center}\rule{0.5\linewidth}{0.5pt}\end{center}

\subsubsection{View 3: Edit Rules (Human-in-the-Loop Counterfactual
Simulation)}\label{view-3-edit-rules-human-in-the-loop-counterfactual-simulation}

The Edit Rules view allows clinical domain experts to inspect, modify,
add, or delete rules, and observe the immediate mathematical impact on
system fidelity.

\begin{verbatim}
+------------------------------------------------------------------------------+
|  EDIT THE RULE BASE                 [ + Add Rule ]  [ Reset to Trained ]     |
|  Holdout Test Fidelity: 92.2% | Conditions: 12 | Rules: 4                    |
+------------------------------------------------------------------------------+
|  RULE 1 (Healthy Baseline)                                                   |
|  IF [Glucose is Low ▾] [×] AND [BMI is Low ▾] [×]     THEN Risk: 0.08 [Save] |
|  [ + Add condition ]                                            [Delete Rule]|
+------------------------------------------------------------------------------+
|  RULE 4 (Overt Diabetes)                                                     |
|  IF [Glucose is High ▾] [×] AND [BMI is High ▾] [×]    THEN Risk: 0.85 [Save] |
|  [ + Add condition ]                                            [Delete Rule]|
+------------------------------------------------------------------------------+
\end{verbatim}

\paragraph{How to Interact with the Rule
Editor}\label{how-to-interact-with-the-rule-editor}

\begin{enumerate}
\def\labelenumi{\arabic{enumi}.}
\tightlist
\item
  \textbf{Change a Condition:} Click the dropdown on any condition pill
  (e.g., change \texttt{Glucose\ is\ High} to
  \texttt{Glucose\ is\ Medium}).
\item
  \textbf{Delete a Condition:} Click the small \texttt{×} icon next to a
  condition to remove it from the rule.
\item
  \textbf{Add a Condition:} Click \texttt{+\ Add\ condition} to
  introduce an additional biomarker constraint to an existing rule.
\item
  \textbf{Delete a Rule:} Click \texttt{Delete\ Rule} to remove an
  entire rule from the knowledge base.
\item
  \textbf{Add a Rule:} Click \texttt{+\ Add\ rule} at the top right to
  define an entirely new clinical scenario.
\item
  \textbf{Instantaneous Re-Fitting (Ridge RLS):}

  \begin{itemize}
  \tightlist
  \item
    The moment any rule edit is made, the Python backend executes Ridge
    Regularized Least Squares across the training cohort in
    \textbf{under 3 milliseconds}.
  \item
    Consequent risk values (\(p_r\)) re-optimize instantly.
  \item
    The top banner immediately updates the \textbf{Test-Set Fidelity}
    score (e.g., showing if deleting a condition caused fidelity to drop
    from \(92.2\%\) to \(89.5\%\)).
  \end{itemize}
\item
  \textbf{Reset to Trained Rules:} Click
  \texttt{Reset\ to\ trained\ rules} at any time to undo all manual
  edits and restore the algorithmically optimized model.
\end{enumerate}

\begin{center}\rule{0.5\linewidth}{0.5pt}\end{center}

\subsubsection{View 4: Compare Runs (Pareto Frontier \& Historical
Benchmarks)}\label{view-4-compare-runs-pareto-frontier-historical-benchmarks}

This view provides comprehensive experimental auditing, allowing
researchers and professors to compare multiple runs across different
hyper-parameter settings.

\begin{verbatim}
+------------------------------------------------------------------------------+
|  PARETO FRONTIER: FIDELITY VS. COMPLEXITY                                    |
|                                                                              |
|   Fidelity (%)                                                               |
|   100% +                                                                     |
|        |               $\bullet$ (Run #1: ANFIS 4 rules, 12 cond, 92.2% fidelity)    |
|    90% +         $\bullet$     $\blacktriangle$ (TREPAN Decision Tree: 12 cond, 84.1% fidelity)     |
|        |                                                                     |
|    80% +   ■ (Direct Fuzzy: 12 cond, 78.4% fidelity)                         |
|        +-------+--------------+--------------+--------------+------------►   |
|                5             10             15             20  Conditions    |
+------------------------------------------------------------------------------+
|  RUNS AUDIT TABLE                                                            |
|  Run ID   | Model Type | Rules | Conditions | Test Fidelity | ROC-AUC | Date |
|  Run #1   | ANFIS (NN) |   4   |     12     |     92.2%     |  0.797  | ...  |
|  Run #2   | Tree-Surr  |   6   |     18     |     85.7%     |  0.781  | ...  |
+------------------------------------------------------------------------------+
\end{verbatim}

\paragraph{Elements of the Compare
View}\label{elements-of-the-compare-view}

\begin{enumerate}
\def\labelenumi{\arabic{enumi}.}
\tightlist
\item
  \textbf{The Pareto Frontier Scatter Chart:}

  \begin{itemize}
  \tightlist
  \item
    \textbf{Y-Axis:} Holdout Test-Set Fidelity (\(\%\) agreement with
    the Neural Network).
  \item
    \textbf{X-Axis:} Model Complexity (Total number of active conditions
    across all rules).
  \item
    \textbf{Optimal Region (Top-Left):} Represents the ideal clinical
    trade-off---maximum fidelity with minimum cognitive complexity.
  \item
    \textbf{Legend Markers:}

    \begin{itemize}
    \tightlist
    \item
      \textbf{Blue Solid Circle (\texttt{$\bullet$}):} Glass-Box ANFIS surrogate
      imitating the neural network.
    \item
      \textbf{Red Solid Circle (\texttt{$\bullet$}):} Control fuzzy model
      trained directly on true labels \(\{0, 1\}\).
    \item
      \textbf{Hollow Circle (\texttt{$\circ$}):} Model state after manual
      clinician edits.
    \item
      \textbf{Green Triangle (\texttt{$\blacktriangle$}):} Decision Tree (TREPAN)
      baseline of equivalent size.
    \end{itemize}
  \end{itemize}
\item
  \textbf{Historical Runs Audit Table:}

  \begin{itemize}
  \tightlist
  \item
    Displays every run executed during the session with its exact
    hyper-parameter configuration.
  \item
    Clicking on any historical row immediately reloads that model into
    memory, allowing you to re-explore its patients in View 2 or inspect
    its rules in View 3.
  \end{itemize}
\end{enumerate}

\begin{center}\rule{0.5\linewidth}{0.5pt}\end{center}

\subsection{7. Frequently Asked Questions \& Professor Viva Defense
Guide}\label{frequently-asked-questions-professor-viva-defense-guide}

\subsubsection{Q1: Is this project strictly within the syllabus of our
course
(EEO1003)?}\label{q1-is-this-project-strictly-within-the-syllabus-of-our-course-eeo1003}

\textbf{Yes, 100\%.} Every algorithmic component directly implements
foundational topics from the course domain: - \textbf{Neural Networks:}
Multi-Layer Perceptrons, backpropagation learning, binary cross-entropy,
and gradient optimization (grounded in Simon Haykin and Laurene
Fausett). - \textbf{Fuzzy Logic:} Linguistic variables, continuous
membership functions, T-norms, zero-order Takagi-Sugeno-Kang fuzzy
inference, and defuzzification (grounded in Timothy J. Ross and S.N.
Sivanandam). - \textbf{Neuro-Fuzzy Synthesis:} ANFIS architecture (Jang,
1993) and Fuzzy C-Means clustering (Bezdek, 1981). - The \emph{only}
modern addition is the overarching problem formulation: using the
neuro-fuzzy system as a pedagogical surrogate for Explainable AI (XAI)
rather than a simple standalone classifier.

\subsubsection{Q2: Why train a Neural Network first if we end up using
Fuzzy
Rules?}\label{q2-why-train-a-neural-network-first-if-we-end-up-using-fuzzy-rules}

Because neural networks possess superior non-linear gradient search
capabilities that find complex decision boundaries in high-dimensional
continuous data. A fuzzy system initialized from scratch often gets
trapped in local optima or suffers from the curse of dimensionality
(\(3^8 = 6{,}561\) rules). By using the \textbf{Pedagogical Surrogate
Philosophy} (Andrews et al., 1995), the neural network acts as an
``oracle teacher'', smoothing out noise in the raw dataset. ANFIS then
distills that teacher's knowledge into 4 clean, human-auditable rules.

\subsubsection{Q3: Why use Zero-Order TSK instead of Mamdani or
First-Order
TSK?}\label{q3-why-use-zero-order-tsk-instead-of-mamdani-or-first-order-tsk}

\begin{itemize}
\tightlist
\item
  \textbf{Against Mamdani:} Mamdani uses fuzzy output sets requiring
  numerical integration (centroid defuzzification), which is
  computationally expensive and difficult to optimize with gradient
  descent.
\item
  \textbf{Against First-Order TSK:} First-order TSK uses linear
  hyperplanes in the THEN-part
  (\(\text{Risk} = a_0 + a_1 x_1 + \dots\)). If the THEN-part is an
  equation, a doctor still cannot interpret the rule at a glance.
\item
  \textbf{Zero-Order TSK Advantage:} The THEN-part is a single,
  intuitive scalar percentage (\(\text{Risk} = 0.85\)), providing
  immediate cognitive clarity while remaining 100\% differentiable for
  hybrid learning.
\end{itemize}

\subsubsection{Q4: Why use the Product T-Norm instead of the Min
operator?}\label{q4-why-use-the-product-t-norm-instead-of-the-min-operator}

The standard Zadeh Min operator (\(\min(\mu_1, \mu_2)\)) has
discontinuous, non-differentiable switching points:
\(\frac{\partial}{\partial \mu_1} \min(\mu_1, \mu_2) = 0\) whenever
\(\mu_1 > \mu_2\). This blocks gradient backpropagation to the
non-minimum feature. The \textbf{Product T-Norm}
(\(w_r = \mu_1 \times \mu_2 \times \dots\)) is smooth, infinitely
differentiable everywhere, and ensures that every condition contributes
proportional gradient information during backward premise tuning.

\subsubsection{Q5: How does the system prevent semantic
collapse?}\label{q5-how-does-the-system-prevent-semantic-collapse}

In unconstrained ANFIS, gradient descent often causes \emph{Low} and
\emph{High} membership curves to invert or collapse into sharp delta
spikes. Our engine enforces \textbf{three strict geometric projection
steps} after every Adam training step: 1.
\(c_{\text{Low}} \le c_{\text{Med}} - 0.05\) and
\(c_{\text{Med}} \le c_{\text{High}} - 0.05\) (ordering constraint). 2.
\(|c - c^{(0)}| \le 0.10\) (drift bounds from data cluster centers). 3.
\(\sigma \ge 0.05\) (width floor preventing delta spikes).

\subsubsection{Q6: What happens if an atypical or anomalous patient is
tested?}\label{q6-what-happens-if-an-atypical-or-anomalous-patient-is-tested}

In black-box models, out-of-distribution patients cause \textbf{silent
catastrophic failures} (the network confidently emits an incorrect risk
score). Glass-Box Insight actively computes the match-strength metric:
\[\mathcal{S}_{\text{match}} = \max_{r=1 \dots R} w_r\] If
\(\mathcal{S}_{\text{match}} < 0.20\), the patient falls outside the
support of all extracted clinical prototypes. The dashboard immediately
triggers an amber/red \textbf{``Weak Match: Out of Distribution''}
diagnostic flag, warning clinicians to disregard automated scoring and
conduct direct manual pathology tests.

\begin{center}\rule{0.5\linewidth}{0.5pt}\end{center}

\emph{Document compiled and verified for the Department of Electrical
Engineering, Punjab Engineering College (Deemed to be University),
Chandigarh.}

\end{document}
